% =====================================================================
\chapter{Телеметрия}
\label{ch:telemetry}
% =====================================================================
\chapterintro
  {как найти датчики, что означают параметры связи ExpressLRS, как сделать вычисляемые
   датчики, тревоги и экраны телеметрии}
  {модель с приёмником, поддерживающим телеметрию; для батареи --- полётный контроллер или
   датчик напряжения}

\section{Откуда берётся телеметрия}

Приёмник измеряет качество радиосвязи и отправляет его обратно на пульт. Если к приёмнику
подключён полётный контроллер (по CRSF для ELRS или по SmartPort/FPort для FrSky), контроллер
добавляет свои данные: напряжение и ток батареи, углы, GPS, режим полёта, высоту.
В EdgeTX каждое значение --- это \textbf{датчик} (\emph{sensor}), который можно использовать
как источник везде: в логических переключателях, голосе, логах, Lua-скриптах.

\section{Поиск датчиков}

\begin{enumerate}
  \item Включите модель, дождитесь связи.
  \item \mpath{MDL\sep Telemetry\sep Discover new sensors} --- EdgeTX начнёт добавлять
        все датчики, которые приходят от приёмника.
  \item Через 5--10 секунд остановите поиск (\menu{Stop discovery}).
\end{enumerate}

Новые датчики появятся списком. Мигающее или «?» значение означает, что данные сейчас
не приходят. Чтобы начать с чистого листа, используйте \menu{Delete all sensors}.

\section{Датчики ExpressLRS}

\begin{center}\small
\renewcommand{\arraystretch}{1.15}
\begin{tabularx}{\linewidth}{@{}p{16mm}X@{}}
\toprule
Датчик & Значение \\
\midrule
\sw{1RSS}, \sw{2RSS} & уровень сигнала на антеннах 1 и 2 приёмника, дБм (чем ближе к 0, тем лучше;
  у предела чувствительности --- около $-105\ldots-115$) \\
\sw{RQly} & качество связи на приёмнике (Link Quality), \% принятых пакетов. \textbf{Главный
  показатель}: ниже 70\,\% --- пора разворачиваться \\
\sw{RSNR} & отношение сигнал/шум на приёмнике, дБ \\
\sw{ANT} & активная антенна приёмника \\
\sw{RFMD} & текущая частота пакетов (режим) \\
\sw{TPWR} & текущая мощность передатчика, мВт \\
\sw{TRSS}, \sw{TQly}, \sw{TSNR} & то же, но на стороне пульта (приём телеметрии) \\
\sw{RxBt} & напряжение батареи (от контроллера или приёмника) \\
\sw{Curr}, \sw{Capa}, \sw{Bat\%} & ток, израсходованная ёмкость, остаток заряда \\
\sw{GPS}, \sw{GSpd}, \sw{Hdg}, \sw{Alt}, \sw{Sats} & координаты, скорость, курс, высота,
  число спутников \\
\sw{Ptch}, \sw{Roll}, \sw{Yaw} & углы ориентации \\
\sw{FM} & режим полётного контроллера (текстом: \code{ACRO}, \code{ANGL}, \code{!FS!}) \\
\bottomrule
\end{tabularx}
\end{center}

\section{Настройка датчика}

\btn{ENTER} на датчике → редактирование:
\begin{itemize}
  \item \menu{Name} --- имя (до 4 символов);
  \item \menu{Unit}, \menu{Precision} --- единица и число знаков после запятой;
  \item \menu{Ratio}, \menu{Offset} --- масштаб и смещение (для аналоговых датчиков);
  \item \menu{Auto Offset}, \menu{Filter} --- обнуление при первом значении, усреднение;
  \item \menu{Persistent} --- хранить последнее значение между включениями (для счётчиков);
  \item \menu{Logs} --- записывать ли датчик в логи.
\end{itemize}

\section{Вычисляемые датчики}

\menu{Add new sensor} → \menu{Type} = \menu{Calculated}. Формулы:

\begin{center}\small
\begin{tabularx}{\linewidth}{@{}p{28mm}X@{}}
\toprule
\menu{Add}, \menu{Average}, \menu{Min}, \menu{Max}, \menu{Multiply} & над несколькими датчиками \\
\menu{Totalize} & интеграл по времени (например, ток → ёмкость) \\
\menu{Cell} & напряжение самой слабой (или выбранной) банки из датчика элементов \\
\menu{Consumption} & израсходованная ёмкость, мА·ч, по датчику тока \\
\menu{Distance} & расстояние от точки взлёта по GPS \\
\bottomrule
\end{tabularx}
\end{center}

\begin{tip}
Напряжение \emph{на банку} удобнее получить на стороне полётного контроллера: в Betaflight
есть параметр \code{report\_cell\_voltage = ON} --- тогда \sw{RxBt} придёт уже в пересчёте
на одну банку, и пороги тревог не надо переделывать для 4S/6S.
\end{tip}

\section{Тревоги связи}

На странице \menu{Telemetry} есть пороги \menu{Low alarm} и \menu{Critical alarm} для
качества связи. Для ELRS ориентируйтесь на \sw{RQly}: предупреждение около 70\,\%, критично ---
около 50\,\%. Пульт будет говорить «Signal low» / «Signal critical». При полной потере связи
прозвучит «Telemetry lost», при восстановлении --- «Telemetry recovered».

\section{Экраны телеметрии}

Кнопка \btn{TELE} на главном экране показывает до четырёх экранов модели. Их содержимое
задаётся в \menu{Telemetry} в разделе \menu{Screens}:

\begin{description}[font=\sffamily\bfseries\color{etx}]
  \item[Nums] сетка из значений (до 4 строк × 3 столбца);
  \item[Bars] полосы с минимумом и максимумом;
  \item[Script] экран, нарисованный Lua-скриптом из папки \code{SCRIPTS/TELEMETRY}
        (глава~\ref{ch:lua}).
\end{description}

\begin{center}
\lcd{\inv{Quad-5in}\hfill\inv{TELE 1/2}\\[2pt]
RxBt\ \ 15.8V\hfill RQly\ \ 100\%\\
Curr\ \ \ 12.3A\hfill 1RSS\ \ -54\\
Capa\ \ \ 812\hfill TPWR\ \ 100mW\\
FM\ \ \ \ \ ACRO\hfill Tmr1\ \ 02:37}
\end{center}

\section{Сброс телеметрии}

Минимумы, максимумы и счётчики (\sw{Capa}) хранятся до сброса. Добавьте специальную функцию
\menu{Reset} → \menu{Telemetry} на переключатель или долгое нажатие кнопки, либо
сбросьте из меню главного экрана (долгое \btn{ENTER} → \menu{Reset}).

\begin{tasks}
\begin{enumerate}
  \item Найдите датчики своей модели и запишите, какие из них пришли.
  \item Настройте экран \menu{Nums} с \sw{RxBt}, \sw{RQly}, \sw{1RSS}, \sw{TPWR}.
  \item Сделайте голосовое предупреждение о низком напряжении (главы~\ref{ch:logic}
        и~\ref{ch:special}).
  \item Отойдите с моделью (без пропеллеров!) на 30--50\,м при минимальной мощности и
        посмотрите, как меняются \sw{1RSS} и \sw{RQly}.
\end{enumerate}
\end{tasks}
